% Source: https://artofproblemsolving.com/community/c3h1785758_sum_of_remainder
% Stable identity: aops-topic-1785758-q01
The remainders of consecutive integers modulo $5$ run cyclically through $0,1,2,3,4$, so every block of $5$ consecutive integers contributes $0+1+2+3+4=10$. Write the number of terms as $k+1=5q+r$ with $0\le r\le4$; then the sum is $10q+s$, where $s$ is the sum of the remainders of the first $r$ terms $n,\dots,n+r-1$, and $0\le s\le1+2+3+4=10$.

Since the sum $2019$ ends in $9$, we need $s=9$ and $q=201$ (if $s$ could be $19$ we would need $r\ge5$). Sums of $r<5$ cyclically consecutive remainders equal to $9$: $r=3$ with remainders $2,3,4$, or $r=4$ with remainders $2,3,4,0$; $r\le2$ gives at most $3+4=7$. In both cases $n\equiv2\pmod5$.

\textbf{Smallest $k$:} $k+1=5\cdot201+3=1008$, so $k=1007$. \textbf{Smallest $n$:} $n\equiv2\pmod 5$ forces $n\ge2$, and $n=2$ works with $k=1007$: the numbers $2,3,\dots,1009$ give $2+3+4$ and then $201$ full cycles, total $9+2010=2019$.

So both minima are attained together: $\boxed{n=2,\ k=1007}$.
